\documentclass{article}
\usepackage{amsmath}
\begin{document}
\section*{A wheel may branch physically without a diagnostic bit}
\paragraph{What changed.}
PRs 844--845 require each issued target to be physically realizable
in both possible worlds before diagnosis, with the same realized
wheel profile. Replace that contract by one allowing specified
wheels to saturate an excessive downward speed request locally.
From actual speed \(v\), brake bound \(b\), and absolute request
\(0\le u\le v\), realized next speed on such a wheel is exactly
\(\max(u,v-b)\). On other wheels a request with \(u<v-b\) is
inadmissible. Even the requested command cannot accelerate in
either world. The responses remain linearly interpolated over
each unit step. This is a new stipulated plant, not a reinterpretation
of the earlier proof and not a claim about the response of Gazebo,
a generic torque ceiling, or hardware.

\paragraph{Paired fault and mission.}
Start with wheel speeds \((4,4)\), radius one and spacing two.
In the straight world the brake bounds are \((1,2)\) and the
terminal heading must be zero; in the turned world they are
\((2,1)\) and the terminal heading must be one radian.
Stop both wheels by step five. Let \(X_L,X_R\) be twice their
distances, so arc length is \((X_L+X_R)/4\) and the scaled terminal
difference is either zero or four. A blind law sends the same
time-indexed absolute requests in both worlds and reads no feedback.
A wheel with weak bound one travels at least eight distance units
before reaching rest; thus its scaled distance is at least 16.

\paragraph{Exact boundary by permitted saturation port.}
If neither wheel may saturate, the two realized paths are identical
under any common admissible request sequence, so the two headings
cannot be attained. If only the right wheel may saturate, left
distance \(X_L\) is common and at least 16 because it is weak in
the straight world. The turned world consequently needs
\(X_L+X_R=2X_L+4\ge36\), hence arc nine.
If the left wheel may saturate, arc eight is attainable, with
or without right saturation; no solution can beat eight because
the straight state itself needs that much. These lower bounds
hold also for real-valued requests in the declared interpolation
model.

For example, the common sequence
\[
 (3,3),(2,1),(0,1),(0,1),(0,0)
\]
is valid when both wheels may saturate. It produces scaled
distances \((16,16)\) and \((14,18)\) in the two worlds, with
the same arc eight and two different correct terminal headings.
The left-only and right-only positive witnesses likewise use
nonincreasing desired speeds at both wheels. The effective branch comes from physical response,
not measurement or an unobserved conditional command. Adding
an exact electronic cap common to both worlds instead would
not supply this branch. A weak-wheel label alone does not say
which contract describes the plant.

\paragraph{Verification and missing obligations.}
A paired-state search and independent factorization through
complete individual-wheel request words agree in 64 scoped instances and 1,760 feasible heading-pair queries. The Gazebo snippets check 30 written target pairs
for six positive profiles; the modeled local response is applied
separately after the native writes. Lean checks the discrete
distance bookkeeping. The exact left/right obstruction is a
physical-model statement, not an advantage over an equally
informed actuator allocator. Tracking, repeatability of the
specific saturation response, tolerance to parameter errors,
swept clearance in obstacles, I1 and independent mission
validation remain open. The original diagnosis cutoff still
applies on its original no-clipping interface.
\end{document}
